% answers_short.tex —— 2009 年·普通高等学校招生考试·天津卷·文科 · 简短答案速查
% 由 tools/make_short_answers.py 自 answers.tex 的 \gksolhead/\gkanswer 提取，
% 并与 2009/2009_数学真题_天津文_答案要点.md 逐题比对后生成，勿手改（改源 MD 后重跑）。

\begin{gknotice}
\textbf{使用说明}：本节只给最终答案便于快速核对，完整推导见解析版。答案为编写者独立推导，\textbf{非官方标准答案}。
\end{gknotice}

\gkpart{选择题}
\begin{center}\begin{tabular}{ccccccccccc}
\toprule
题号 & 1 & 2 & 3 & 4 & 5 & 6 & 7 & 8 & 9 & 10 \\
\midrule
答案 & D & B & A & C & B & C & D & A & C & A \\
\bottomrule
\end{tabular}\end{center}

\gkpart{填空题}
\begin{itemize}[leftmargin=2.6em]
  \item[\textbf{11.}] $2$。
  \item[\textbf{12.}] $\sqrt3$。
  \item[\textbf{13.}] $\{2,4,6,8\}$。
  \item[\textbf{14.}] $1$。
  \item[\textbf{15.}] $-2$。
  \item[\textbf{16.}] $\left(\dfrac{25}{9},\dfrac{49}{16}\right]$。
\end{itemize}

\gkpart{解答题}
\begin{itemize}[leftmargin=2.6em]
  \item[\textbf{17.}] （1）$AB=2\sqrt5$；（2）$\dfrac{\sqrt2}{10}$。
  \item[\textbf{18.}] （1）$2,3,2$；（2）$\dfrac{11}{21}$。
  \item[\textbf{19.}] （1）（2）证明见解析；（3）$\dfrac13$。
  \item[\textbf{20.}] （1）$a_n=4n-3$；（2）$q=-2$；（3）证明见解析。
  \item[\textbf{21.}] （1）$1$；（2）增区间 $(1-m,1+m)$，减区间 $(-\infty,1-m)$、$(1+m,+\infty)$，极大值 $\dfrac{(2m-1)(m+1)^2}{3}$，极小值 $-\dfrac{(2m+1)(m-1)^2}{3}$；（3）$\left(\dfrac12,\dfrac{\sqrt3}{3}\right)$。
  \item[\textbf{22.}] （1）$\dfrac{\sqrt3}{3}$；（2）$\pm\dfrac{\sqrt2}{3}$；（3）$\pm\dfrac{2\sqrt2}{5}$。
\end{itemize}
